\section{Soil Dataset Description}

The soil dataset used in this study consists of \textbf{23,262 grid cells}
($0.1^{\circ}$ resolution) covering Algeria and Tunisia . Each row corresponds to a
regular grid-cell centroid (\texttt{x}, \texttt{y}) together with topsoil
(0--20\,cm) physical and chemical attributes extracted from the Harmonized World
Soil Database v2 (HWSD2). The total memory usage for the soil data is approximately \textbf{5.90 MB}. All features originate either from the D1 layer
(topsoil) or categorical summaries at the Soil Mapping Unit (SMU) level.

\subsection*{Spatial Columns}
\textbf{x, y} (\textit{longitude, latitude}) specify the centroid of each cell.
These coordinates enable spatial alignment with land cover, climate and
ignition datasets used in fire modelling \citep{Giglio2013}.

\subsection*{Identification Key}
\textbf{HWSD2\_SMU\_ID} uniquely identifies the soil mapping unit associated with
each cell. This ID links the grid to HWSD2 pedological information and ensures
consistent reproducibility \citep{FAO2021}.

\subsection*{Texture and Particle-Size Distribution}
\textbf{COARSE} (\,\%) represents the proportion of fragments $>2$\,mm, influencing
water retention and drying rates. High coarse content increases surface dryness
and ignition probability \citep{Prieto2019}.  
\textbf{SAND}, \textbf{SILT} and \textbf{CLAY} (\,\%) describe the granulometric
composition. Sandy soils drain rapidly and dry faster, whereas clay-rich soils
retain moisture and generally reduce fire susceptibility
\citep{Bowman2009,Saatchi2012}.  
\textbf{TEXTURE\_USDA} and \textbf{TEXTURE\_SOTER} provide categorical texture
classes (USDA and SOTER standards). Texture strongly regulates infiltration,
soil moisture availability and vegetation composition \citep{USDA2017,Batjes2016}.

\subsection*{Density and Structural Properties}
\textbf{BULK} (g\,cm$^{-3}$) measures dry soil density. Higher bulk density limits
root penetration and reduces water storage capacity \citep{Certini2005}.  
\textbf{REF\_BULK} is the expected bulk density estimated from pedotransfer
functions, used to detect compaction anomalies \citep{Reynolds2018}.

\subsection*{Organic Matter and Nutrients}
\textbf{ORG\_CARBON} (\%) represents organic carbon content, a driver of both
soil moisture retention and surface fuel accumulation \citep{DeGroot2009}.  
\textbf{TOTAL\_N} (\%) influences vegetation productivity and post-fire biomass
recovery \citep{Pausas2012}.  
\textbf{CN\_RATIO} (C:N) governs decomposition rates; high ratios slow down decay
and preserve surface litter, increasing flammability \citep{Fernandez2018}.

\subsection*{Chemical and Exchange Properties}
\textbf{PH\_WATER} indicates soil acidity. Extreme pH affects nutrient cycling
and decomposition, altering fuel loads \citep{Certini2005b}.  
\textbf{CEC\_SOIL}, \textbf{CEC\_CLAY} and \textbf{CEC\_EFF} (cmol\,kg$^{-1}$)
represent nutrient-holding capacity. High CEC improves water retention and thus
reduces fire proneness \citep{Brady2017,Jenny1941}.  
\textbf{TEB} (Total Exchangeable Bases) and \textbf{BSAT} (\%) describe soil
fertility; productive soils often accumulate more biomass, increasing available
fuel \citep{Archibald2013}.  
\textbf{ALUM\_SAT} (\%) and \textbf{ESP} (\%) capture aluminum and sodium
saturation. High ESP causes soil crusting and sparse vegetation, producing
patchy fire propagation patterns \citep{Qadir2014}.  
\textbf{TCARBON\_EQ} (total carbon equivalents) and \textbf{GYPSUM} (\%) indicate
soil development stage in arid systems. Gypsiferous and carbonate-rich soils
often host low vegetation cover and reduced fire intensity \citep{Herrero2009}.  
\textbf{ELEC\_COND} (dS\,m$^{-1}$) measures salinity; high salinity stresses
vegetation and correlates with lower fuel loads \citep{RuizSinoga2012}.

Overall, these variables collectively describe the soil moisture regime, nutrient
status, vegetation-support capacity and physical structure of the topsoil, all
of which are key environmental determinants of wildfire ignition and spread.
